% Figure: two-panel residual diagnostic for the end-to-end case study.
% Left: residuals against fitted values on the log-life scale, showing a
% constant-width band consistent with the homoscedastic model. Right: a
% normal quantile-quantile plot of the standardised residuals, close to
% the reference line with mild departure in the upper tail. Coordinates
% are representative residuals from the worked fit in the text.
\begin{figure}[htbp]
\centering
\begin{tikzpicture}
% Left panel: residuals vs fitted
\begin{axis}[
    name=resid,
    width=7.2cm, height=5.4cm,
    xmin=3.6, xmax=5.2,
    ymin=-0.45, ymax=0.45,
    axis lines=left,
    xlabel={Fitted $\widehat{\log_{10} N}$},
    ylabel={Residual},
    title={\small Residuals vs fitted},
    ytick={-0.4,-0.2,0,0.2,0.4},
    grid=both,
    grid style={dashed, draw=enggray!20},
]
\addplot[only marks, mark=*, mark size=1.3pt, engblue] coordinates {
  (3.72,0.06) (3.85,-0.11) (3.98,0.14) (4.05,-0.05) (4.18,0.09)
  (4.22,-0.18) (4.35,0.03) (4.41,0.21) (4.48,-0.09) (4.55,0.12)
  (4.62,-0.14) (4.70,0.07) (4.78,-0.04) (4.85,0.17) (4.92,-0.10)
  (5.01,0.02) (4.10,-0.22) (4.30,0.10) (4.66,-0.06) (4.88,0.30)
};
\addplot[engred, thick] coordinates {(3.6,0) (5.2,0)};
\end{axis}
% Right panel: normal QQ plot
\begin{axis}[
    at={(resid.east)}, anchor=west, xshift=1.1cm,
    width=7.2cm, height=5.4cm,
    xmin=-2.3, xmax=2.3,
    ymin=-2.3, ymax=2.6,
    axis lines=left,
    xlabel={Theoretical normal quantile},
    ylabel={Standardised residual},
    title={\small Normal Q--Q},
    grid=both,
    grid style={dashed, draw=enggray!20},
]
\addplot[engred, thick, domain=-2.2:2.2] {x};
\addplot[only marks, mark=*, mark size=1.3pt, engblue] coordinates {
  (-1.87,-1.74) (-1.38,-1.30) (-1.09,-1.15) (-0.85,-0.79)
  (-0.64,-0.71) (-0.45,-0.40) (-0.28,-0.31) (-0.12,-0.05)
  (0.04,0.10) (0.20,0.18) (0.37,0.33) (0.54,0.62) (0.73,0.70)
  (0.95,0.99) (1.21,1.28) (1.53,1.41) (2.05,2.34)
};
\end{axis}
\end{tikzpicture}
\caption[Residual diagnostics for the case-study fit]{Residual
diagnostics for the case-study regression of fatigue life on stress
amplitude. The left panel plots residuals against fitted values on
the $\log_{10}$-cycles scale; the band is of roughly constant width
about zero, consistent with the homoscedastic Gaussian-error model
the analysis assumed after the log transform. The right panel is a
normal quantile-quantile plot of the standardised residuals; the
points track the reference line through the bulk of the range with a
single upper-tail point sitting above the line. The mild upper-tail
departure is the kind of finding the engineer reports rather than
suppresses: it weakens the Gaussian-tail assumption slightly without
overturning the fit, and it is the reason the case study reports a
bootstrap interval alongside the parametric one.}
\label{fig:vol02:ch14:casestudy-diagnostics}
\end{figure}
